Generalist-robot safety must include \emph{initiation
authorization}: deciding whether a robot may take its first
hard-to-undo social action when the situation is still unclear. An
engagement score measures how receptive someone seems. It does not
answer whether the robot should act when that signal is ambiguous
\citep{rich2010,broek2024proactive,shi2011spatial}.

Deployed stacks often collapse two distinct safety decisions into
one pipeline. \emph{Physical safety} limits motion---collisions,
force, joint limits \citep{zhang2025safevla}.
\emph{Social-interaction safety} governs how the robot behaves once
interaction has already started, including recovery when a first act
misfires \citep{akalin2021,tian2021socialerrors,harland2025aiapology}.
We argue for a third layer that this picture largely omits:
\emph{initiation authorization}---whether the robot may start at all.
Standard designs rarely treat it as a separate decision. The harm
often looks like an
ordinary first move: a mistimed greeting that interrupts a call
\citep{edirisinghe2024shopworker}, or an end-to-end VLA that reaches
for an object before the person has agreed
\citep{openvla-2024}. A jailbroken planner can
authorize a harmful first move as readily as a harmful first word
\citep{robey2024jailbreaking}.

Most deployed stacks treat initiation as a fixed threshold:
$\textsc{act/speak} \Leftrightarrow e_t > \theta$
\citep{bohus2014}. That treats an engagement score, or a single VLA
rollout, as permission to act, with no context, no user preference,
and no evidence from low-cost probes
\citep{skantze2014speech,bohus2014}.
Figure~\ref{fig:initiation} contrasts this passive threshold with
\textsc{pas} (probe--authorize--speak): small reversible probes
before any hard-to-undo action, an explicit bold/conservative setting,
and runtime checks before foundation-model dialogue.

% Figure 1 — compact two-column teaser
\begin{figure*}[t]
\centering
\resizebox{0.98\textwidth}{!}{%
\begin{tikzpicture}[
  font=\scriptsize,
  >=Stealth,
  arr/.style={-{Stealth[length=2mm,width=1.6mm]}, line width=0.9pt},
  panel/.style={rounded corners=3pt, line width=0.7pt},
  hdr/.style={font=\footnotesize\bfseries, text=white},
  card/.style={rounded corners=2pt, align=center, font=\tiny,
    inner sep=3pt, minimum width=1.35cm, minimum height=0.7cm},
  mod/.style={rounded corners=2pt, align=center, font=\tiny,
    inner sep=4pt, line width=0.6pt, text width=2.5cm},
]
\def\colW{3.55}
\def\yBot{0.35}
\def\yTop{3.55}
\def\yHdrBot{3.15}
\def\xA{0.05}
\def\xB{3.75}
\def\xC{7.45}

\fill[MenSlate, rounded corners=2pt] (0,3.72) rectangle (11.05,4.02);
\node[font=\footnotesize\bfseries, text=white] at (5.52,3.87)
  {Initiation authorization under uncertainty};

% Col 1
\fill[MenSlateLight, panel] (\xA,\yBot) rectangle (\xA+\colW,\yTop);
\draw[MenSlate!45, panel] (\xA,\yBot) rectangle (\xA+\colW,\yTop);
\fill[MenSlate!75, rounded corners=2pt] (\xA+0.1,\yHdrBot) rectangle (\xA+\colW-0.1,\yTop-0.06);
\node[hdr, anchor=west] at (\xA+0.2,\yTop-0.22) {Social context};
\node[card, draw=MenCoral!70, fill=MenCoralLight!40] at (\xA+0.95,2.55) {on phone\\busy};
\node[card, draw=MenSlate!50, fill=white] at (\xA+2.55,2.55) {glance\\unclear};
\node[card, draw=MenTeal!70, fill=MenTealLight!35] at (\xA+0.95,1.75) {gaze\\probe OK};
\node[card, draw=MenCoral!70, fill=MenCoralLight!40] at (\xA+2.55,1.75) {wrong\\addressee};
\fill[MenCoralLight!55, rounded corners=2pt] (\xA+0.15,0.55) rectangle (\xA+\colW-0.15,1.35);
\node[font=\tiny\bfseries, text=MenCoral!90!black, anchor=west] at (\xA+0.25,1.18) {Passive threshold};
\node[font=\tiny, text=MenSlate, anchor=west] at (\xA+0.25,0.82) {$\texttt{SPEAK} \Leftrightarrow e_t>\theta$};

% Col 2
\fill[MenTealLight!30, panel] (\xB,\yBot) rectangle (\xB+\colW,\yTop);
\draw[MenTeal!50, panel] (\xB,\yBot) rectangle (\xB+\colW,\yTop);
\fill[MenTeal!85, rounded corners=2pt] (\xB+0.1,\yHdrBot) rectangle (\xB+\colW-0.1,\yTop-0.06);
\node[hdr, anchor=west] at (\xB+0.2,\yTop-0.22) {Safety-by-practice};
\node[mod, draw=MenTeal!65, fill=white] (m1) at (\xB+1.78,2.55) {\textbf{Probes}\\never \texttt{speak}};
\node[mod, draw=MenSlate!40, fill=MenSlateLight!50] (m2) at (\xB+1.78,1.85) {\textbf{Evidence}\\receptivity};
\node[mod, draw=MenGold!75!black, fill=MenGold!6] (m3) at (\xB+1.78,1.15) {\textbf{Auth.\ gate}\\sole \texttt{speak}};
\draw[arr, MenTeal!70] (m1) -- (m2); \draw[arr, MenTeal!70] (m2) -- (m3);

% Col 3
\fill[white, panel] (\xC,\yBot) rectangle (\xC+\colW,\yTop);
\draw[MenGold!55, panel] (\xC,\yBot) rectangle (\xC+\colW,\yTop);
\fill[MenGold!80!black, rounded corners=2pt] (\xC+0.1,\yHdrBot) rectangle (\xC+\colW-0.1,\yTop-0.06);
\node[hdr, anchor=west] at (\xC+0.2,\yTop-0.22) {Outcomes};
\fill[MenCoralLight!40, rounded corners=2pt] (\xC+0.15,1.95) rectangle (\xC+\colW-0.15,3.05);
\node[font=\tiny\bfseries, text=MenCoral!90!black, anchor=west] at (\xC+0.25,2.88) {Passive};
\node[font=\tiny] at (\xC+1.78,2.45) {``Hi!'' $\times$ awkward};
\fill[MenTealLight!45, rounded corners=2pt] (\xC+0.15,0.55) rectangle (\xC+\colW-0.15,1.85);
\node[font=\tiny\bfseries, text=MenTeal!90!black, anchor=west] at (\xC+0.25,1.68) {Authorized};
\node[font=\tiny] at (\xC+1.78,1.15) {probe $\rightarrow$ ``May I help?'' $\checkmark$};

\draw[arr, MenSlate!50] (\xA+\colW+0.02,2.05) -- (\xB-0.02,2.05)
  node[midway, above=1pt, font=\tiny, text=MenSlate] {evidence};
\draw[arr, MenTeal!65] (\xB+\colW+0.02,2.05) -- (\xC-0.02,2.05)
  node[midway, above=1pt, font=\tiny, text=MenTeal] {policy};

\fill[MenSlateLight, rounded corners=2pt, draw=MenSlate!18, line width=0.5pt]
  (0.05,-0.08) rectangle (11.0,0.30);
\node[font=\tiny, align=center, inner sep=0pt] at (5.52,0.11) {%
  \textbf{Generalist safety} needs context and preference---harm begins at the \textbf{first irreversible commitment}.%
};
\end{tikzpicture}%
}
\caption{Initiation authorization under uncertainty (left to right).
\textbf{Social context} gives ambiguous cues---a busy phone call, an unclear
glance, a gaze suitable for probing, or a wrong addressee---while many
deployed stacks still use a \emph{passive threshold}
($\texttt{speak} \Leftrightarrow e_t>\theta$) that treats an engagement score
as permission to speak.
\textbf{Safety-by-practice} (\textsc{pas}) instead runs reversible
\emph{probes} that never emit \texttt{speak}, accumulates receptivity
\emph{evidence}, and opens an authorization \emph{gate} as the sole path to
first speech.
\textbf{Outcomes} contrast a passive premature greeting (``Hi!'' $\rightarrow$
awkward) with an authorized sequence (probe $\rightarrow$ ``May I help?''
$\rightarrow$ acceptable).
Arrows mark how context supplies evidence and how policy selects the outcome.
The figure's point is not collision avoidance: a collision-free robot can still
harm a person with its first social commitment.}
\label{fig:initiation}
\end{figure*}


Building on this distinction, we make three contributions:
\begin{enumerate}
\item We introduce \emph{initiation authorization} as a distinct safety
layer, separate from physical shields, post-plan VLA guardrails,
post-dialogue preference alignment, and engagement optimization.
\item We illustrate \textsc{pas} (probe--authorize--speak) on a physical
robot and propose an evaluation protocol that compares three
initiation policies and states how the framing could fail.
\item We elaborate on future work, pose open questions for the community,
and offer provisional statements on metrics, deployment logging,
governance, and where initiation ends and foundation-model generation
begins.
\end{enumerate}
